
\section{Compte rendu du travail en équipe}

Le projet a été réalisé en groupe, ce qui a nécessité une organisation rigoureuse et une communication régulière entre les membres. Le travail en équipe a été globalement bénéfique, mais a également mis en évidence certaines limites.

\subsection{Organisation et répartition du travail}

Nous avons répartis le travail en fonction des modules à réaliser pour le projet.

Jusqu'à la release 1, nous étions organisé de la manière suivante :
\begin{itemize}
    \item Félix : interface graphique
    \item Romain : interprétation JajaCode et compilation MiniJaja vers JajaCode
    \item Lucas : interprétation JajaCode et compilation MiniJaja vers JajaCode
    \item Léo : interprétation MiniJaja, analyse lexicale et syntaxique
    \item Ahmed : interprétation MiniJaja, analyse lexicale et syntaxique
    \item Théo : pile et table des symboles (Mémoire)
    \item Javad : pile et table des symboles (Mémoire)
\end{itemize}
\vspace{0.5cm}
Mais nous avons accumulé du retard sur la partie mémoire vers la fin du sprint 3, donc après la 1ère release nous avons décidé de réorganiser les groupes
en passant Javad à l'interface graphique pour travailler avec félix.



\subsection{Revues de code et communication}

Des revues de code ont été mises en place au travers des merge requests via GitLab afin d’améliorer la qualité globale du projet et de limiter les erreurs d’intégration.
Ces revues ont permis de détecter des problèmes avant la fusion dans la branche de développement, pour garantir un projet fonctionnel sur cette branche.
\vspace{0.5cm}

Nous avons principalement utilisés les merge requests de Gitlab ainsi que discord pour communiquer.
Nous avons plusieurs fois organisés des réunions sur discord afin de discuter de choix de conceptions, ou encore de rétrospectives de sprint.
Les séances de TP ainsi que les TD prévus pour le projet ont également servi de points de synchronisation importants.


\subsection{Avantages et inconvénients du travail en équipe}

Le principal avantage du travail en équipe a été la possibilité de développer plusieurs parties du projet en parallèle, tout en confrontant régulièrement les choix techniques.
Nous avons aussi pu faire du pair programming à plusieurs reprises, comme pour l'implémentation du visiteur MiniJaja.

\vspace{0.5cm}
En revanche, certaines difficultés sont apparues, notamment lors de l’intégration des modules ou en raison de différences de rythme de travail entre les membres.
Pour résoudre ces problèmes, nous avons confrontés nos points de vue lors de réunions et avons ajusté notre organisation en conséquence.
Exemple : lors de la gestion des types pour l'implémentation des noeuds de l'AST, nous avons eu des difficultés à trouver un accord commun.
Etant donné que ce choix a des répercussions sur plusieurs modules, il était important que tout le monde soit d'accord sur la solution à adopter.
Cela nous a alors pris presque une semaine pour trouver une solution et l'implémenter sous la forme d'une énumération dans les modules concernés.